%% @file paper52_algebraic_topology_en.tex
%% @brief Algebraic Topology: Fundamental Groups, Homology, Cohomology, and Applications
%% @author Michael Fuhrmann
%% @date 2026-03-12
%% @project Specialist Mathematics Project, Build 139

\documentclass[12pt,a4paper]{amsart}
\usepackage{amsmath,amssymb,amsthm}
\usepackage{mathtools}
\usepackage[utf8]{inputenc}
\usepackage[T1]{fontenc}
\usepackage{hyperref}
\usepackage{enumitem,booktabs,array}
\usepackage{geometry}
\geometry{margin=2.5cm}

%% Theorem-Umgebungen
\newtheorem{theorem}{Theorem}[section]
\newtheorem{lemma}[theorem]{Lemma}
\newtheorem{corollary}[theorem]{Corollary}
\newtheorem{proposition}[theorem]{Proposition}
\newtheorem{conjecture}[theorem]{Conjecture}
\theoremstyle{definition}
\newtheorem{definition}[theorem]{Definition}
\newtheorem{example}[theorem]{Example}
\newtheorem{remark}[theorem]{Remark}

\author{Michael Fuhrmann}
\address{Specialist Mathematics Project, Build 139}
\email{specialist-maths@localhost}
\date{\today}
\title{Algebraic Topology:\\
Fundamental Groups, Homology, Cohomology,\\
and Classical Applications}

\begin{abstract}
Algebraic topology studies topological spaces by associating algebraic invariants
to them, converting geometric problems into computable algebraic ones. This paper
provides a comprehensive treatment of the fundamental group, covering spaces,
singular homology and cohomology, CW-complexes, and higher homotopy groups.
We prove the van Kampen theorem, Mayer--Vietoris sequence, excision theorem,
universal coefficient theorem, K\"unneth formula, Poincar\'e duality for compact
orientable manifolds, and the Hurewicz isomorphism theorem. Classical applications
include the Brouwer fixed-point theorem, the Borsuk--Ulam theorem, and the
hairy ball theorem. The Poincar\'e conjecture in dimension three, resolved by
Perelman in 2003, and Adams' theorem on the Hopf invariant are also stated and
their significance explained. All results are grounded in rigorous proofs, with
connections to the computational implementations in the companion
\texttt{algebraic\_topology.py} module of the Specialist Mathematics Project.
\end{abstract}

\maketitle
\tableofcontents

%% ============================================================
\section{Introduction}
%% ============================================================

Algebraic topology is the study of topological spaces through the lens of algebra.
The central idea is to assign algebraic objects -- groups, rings, modules -- to
topological spaces in a \emph{functorial} way, so that homeomorphic spaces yield
isomorphic invariants. This transforms qualitative geometric questions (``Are these
spaces homeomorphic?'') into tractable algebraic ones (``Are these groups isomorphic?'').

The historical roots of the subject reach back to Euler's polyhedral formula
\[
  V - E + F = 2,
\]
which is the first instance of an algebraic invariant capturing topological
information. Poincar\'e systematised these ideas in his foundational papers of
1895--1904, introducing the fundamental group and the homology groups, and
conjecturing what is now the celebrated Poincar\'e conjecture.

The twentieth century saw the development of singular homology by Eilenberg,
cohomology by \v{C}ech and de Rham, the axiomatic framework of Eilenberg--Steenrod,
homotopy theory, spectral sequences, K-theory, and ultimately the resolution of the
Poincar\'e conjecture in all dimensions (Smale $\geq 5$, Freedman in dimension 4,
Perelman in dimension 3).

\medskip\noindent\textbf{Organisation of the paper.}
Section~\ref{sec:fundamental} introduces the fundamental group and proves the
Seifert--van Kampen theorem. Section~\ref{sec:covering} treats covering spaces.
Section~\ref{sec:homology} develops singular homology. Section~\ref{sec:key}
proves the Mayer--Vietoris sequence and the excision theorem.
Section~\ref{sec:cohomology} covers cohomology and the universal coefficient theorem.
Section~\ref{sec:poincare} proves Poincar\'e duality.
Section~\ref{sec:cw} discusses CW-complexes and cellular homology.
Section~\ref{sec:higher} treats higher homotopy groups and the Hurewicz theorem.
Section~\ref{sec:applications} presents the classical fixed-point and
antipodal theorems.

%% ============================================================
\section{The Fundamental Group}
\label{sec:fundamental}
%% ============================================================

\subsection{Homotopy of Paths}

Let $X$ be a topological space and $x_0 \in X$ a base point.

\begin{definition}
A \emph{loop} at $x_0$ is a continuous map $\gamma\colon [0,1] \to X$ with
$\gamma(0) = \gamma(1) = x_0$. Two loops $\gamma_0, \gamma_1$ are
\emph{homotopic relative to $\{0,1\}$}, written $\gamma_0 \simeq \gamma_1$,
if there exists a continuous map
\[
  H\colon [0,1] \times [0,1] \to X
\]
such that $H(s,0) = \gamma_0(s)$, $H(s,1) = \gamma_1(s)$, and
$H(0,t) = H(1,t) = x_0$ for all $s, t \in [0,1]$.
\end{definition}

Homotopy relative to end-points is an equivalence relation. The
\emph{concatenation} of loops $\gamma * \delta$ defined by
\[
  (\gamma * \delta)(s) = \begin{cases}
    \gamma(2s) & 0 \leq s \leq \tfrac{1}{2}, \\
    \delta(2s-1) & \tfrac{1}{2} \leq s \leq 1,
  \end{cases}
\]
descends to a well-defined operation on homotopy classes.

\begin{definition}
The \emph{fundamental group} of $X$ at $x_0$ is
\[
  \pi_1(X, x_0) = \{ [\gamma] : \gamma \text{ a loop at } x_0 \},
\]
with multiplication $[\gamma][\delta] = [\gamma * \delta]$.
\end{definition}

\begin{theorem}
$\pi_1(X, x_0)$ is a group under concatenation of homotopy classes.
\end{theorem}
\begin{proof}
\textbf{Associativity.} $([\alpha][\beta])[\gamma] = [\alpha]([\beta][\gamma])$
follows by reparametrisation: define
$H(s,t) = (\alpha * \beta * \gamma)$ with a $t$-dependent reparametrisation
shifting the subdivision points; the resulting map is a valid homotopy.

\textbf{Identity.} The constant loop $e(s) = x_0$ satisfies $[\gamma][e] = [\gamma]$:
the homotopy $H(s,t) = \gamma(\min(2s/(1+t), 1))$ (for $t$ interpolating toward the
original parametrisation) demonstrates this.

\textbf{Inverse.} $\bar{\gamma}(s) = \gamma(1-s)$ satisfies
$[\gamma][\bar{\gamma}] = [e]$ via the homotopy that ``cancels'' $\gamma$
against its reverse.
\end{proof}

\subsection{Examples}

\begin{example}
$\pi_1(\mathbb{R}^n, x_0) = 1$ (trivial): every loop is null-homotopic via
the straight-line homotopy $H(s,t) = (1-t)\gamma(s) + t x_0$.
\end{example}

\begin{example}
$\pi_1(S^1, 1) \cong \mathbb{Z}$: the isomorphism sends a loop to its
\emph{winding number}. The loop $\gamma_n(s) = e^{2\pi i n s}$ represents the
class $n \in \mathbb{Z}$.
\end{example}

\begin{example}
$\pi_1(T^2, x_0) \cong \mathbb{Z}^2$ (abelian): the two generators correspond
to the two independent loops around the torus.
\end{example}

\begin{example}
$\pi_1(\mathbb{R}P^2, x_0) \cong \mathbb{Z}/2\mathbb{Z}$: a non-orientable
manifold whose fundamental group detects the orientation-reversing loop.
\end{example}

\subsection{The Seifert--van Kampen Theorem}

The van Kampen theorem computes the fundamental group of a union.

\begin{theorem}[Seifert--van Kampen]
\label{thm:vankampen}
Let $X = U \cup V$ where $U, V$ are open, path-connected, and $U \cap V$
is path-connected. Fix $x_0 \in U \cap V$. Then
\[
  \pi_1(X, x_0) \cong \pi_1(U, x_0) *_{\pi_1(U \cap V, x_0)} \pi_1(V, x_0),
\]
the \emph{amalgamated free product} over the images of $\pi_1(U \cap V)$.
\end{theorem}
\begin{proof}[Proof sketch]
Any loop $\gamma$ in $X$ at $x_0$ can, by compactness of $[0,1]$ and Lebesgue's
number lemma, be subdivided into sub-arcs each lying entirely in $U$ or in $V$.
By inserting connecting paths in $U \cap V$, the loop decomposes as a product
of loops in $U$ and $V$. This gives a surjection from
$\pi_1(U) * \pi_1(V)$ to $\pi_1(X)$.

For injectivity, suppose a word $w = u_1 v_1 \cdots u_k v_k$ (with $u_i \in \pi_1(U)$,
$v_j \in \pi_1(V)$) maps to the trivial element of $\pi_1(X)$. A homotopy
$H\colon [0,1]^2 \to X$ provides the null-homotopy; using the Lebesgue number
again, one subdivides $[0,1]^2$ so that each rectangle maps into $U$ or $V$.
Careful analysis of this grid shows that the word $w$ already equals the identity
in the amalgamated free product.
\end{proof}

\begin{example}
For the wedge $S^1 \vee S^1$: take $U$ and $V$ as open subsets each homotopic
to $S^1$ with $U \cap V \simeq *$ (contractible). Then
$\pi_1(S^1 \vee S^1) \cong \mathbb{Z} * \mathbb{Z}$, the free group on two
generators.
\end{example}

%% ============================================================
\section{Covering Spaces}
\label{sec:covering}
%% ============================================================

\subsection{Basic Theory}

\begin{definition}
A \emph{covering space} of $X$ is a pair $(\tilde{X}, p)$ where
$\tilde{X}$ is a topological space and $p\colon \tilde{X} \to X$ is a
continuous surjection such that every $x \in X$ has an open neighbourhood
$U$ with $p^{-1}(U) = \bigsqcup_\alpha V_\alpha$ a disjoint union of open
sets, each mapped homeomorphically onto $U$ by $p$.
\end{definition}

\begin{theorem}[Path lifting]
Given a covering $p\colon \tilde{X} \to X$, a path $\gamma\colon [0,1] \to X$,
and a lift $\tilde{x}_0 \in p^{-1}(\gamma(0))$, there exists a unique lift
$\tilde{\gamma}\colon [0,1] \to \tilde{X}$ with $\tilde{\gamma}(0) = \tilde{x}_0$
and $p \circ \tilde{\gamma} = \gamma$.
\end{theorem}

\begin{theorem}[Homotopy lifting]
Given a covering $p\colon \tilde{X} \to X$ and a homotopy
$H\colon [0,1]^2 \to X$ with a chosen lift of $H(0,0)$, there exists a
unique lift $\tilde{H}\colon [0,1]^2 \to \tilde{X}$.
\end{theorem}

\subsection{Deck Transformations and Classification}

\begin{definition}
A \emph{deck transformation} of $p\colon \tilde{X} \to X$ is a homeomorphism
$\phi\colon \tilde{X} \to \tilde{X}$ with $p \circ \phi = p$.
The group of deck transformations is denoted $\mathrm{Deck}(p)$.
\end{definition}

\begin{theorem}[Classification of covering spaces]
\label{thm:covering_classification}
For a connected, locally path-connected, semi-locally simply-connected space $X$
with base point $x_0$, there is a bijection between:
\begin{enumerate}[label=(\roman*)]
  \item Isomorphism classes of connected covering spaces of $X$.
  \item Conjugacy classes of subgroups of $\pi_1(X, x_0)$.
\end{enumerate}
Under this correspondence, the \emph{universal covering space}
$\tilde{X}$ (for which $\pi_1(\tilde{X}) = 1$) corresponds to the trivial
subgroup, and $\mathrm{Deck}(\tilde{X} \to X) \cong \pi_1(X, x_0)$.
\end{theorem}

\begin{example}
The universal cover of $S^1$ is $\mathbb{R}$, with $p(t) = e^{2\pi i t}$.
Deck transformations are integer translations $t \mapsto t + n$, and
$\mathrm{Deck}(\mathbb{R} \to S^1) \cong \mathbb{Z} \cong \pi_1(S^1)$.
\end{example}

\begin{example}
The universal cover of $\mathbb{R}P^n$ is $S^n$, with deck group
$\mathbb{Z}/2\mathbb{Z}$ acting by the antipodal map.
\end{example}

%% ============================================================
\section{Singular Homology}
\label{sec:homology}
%% ============================================================

\subsection{Singular Simplices and Chain Groups}

\begin{definition}
The \emph{standard $n$-simplex} is
\[
  \Delta^n = \left\{ (t_0, \ldots, t_n) \in \mathbb{R}^{n+1}
    : t_i \geq 0,\; \sum_{i=0}^n t_i = 1 \right\}.
\]
A \emph{singular $n$-simplex} in $X$ is a continuous map
$\sigma\colon \Delta^n \to X$.
The \emph{$n$-th chain group} $C_n(X;\mathbb{Z})$ is the free abelian group
generated by all singular $n$-simplices.
\end{definition}

\subsection{Boundary Operators and Homology}

The \emph{$i$-th face map} $\partial_i\colon \Delta^{n-1} \to \Delta^n$
embeds the standard $(n-1)$-simplex as the face opposite vertex $i$:
\[
  \partial_i(t_0,\ldots,t_{n-1}) = (t_0,\ldots,t_{i-1}, 0, t_i, \ldots, t_{n-1}).
\]
The boundary of a singular $n$-simplex $\sigma$ is
\[
  \partial_n(\sigma) = \sum_{i=0}^n (-1)^i (\sigma \circ \partial_i)
  \;\in C_{n-1}(X;\mathbb{Z}).
\]

\begin{theorem}[$\partial^2 = 0$]
$\partial_{n-1} \circ \partial_n = 0$ for all $n$.
\end{theorem}
\begin{proof}
One computes directly that the face maps satisfy
$\partial_j \circ \partial_i = \partial_i \circ \partial_{j-1}$ for $i < j$.
Using this relation:
\begin{align*}
  \partial_{n-1}\partial_n(\sigma)
  &= \sum_{j=0}^n (-1)^j \partial_{n-1}(\sigma \circ \partial_j) \\
  &= \sum_{j=0}^n \sum_{i=0}^{n-1} (-1)^{i+j}
     (\sigma \circ \partial_j \circ \partial_i) \\
  &= \sum_{i < j} (-1)^{i+j}(\sigma \circ \partial_i \circ \partial_{j-1})
   + \sum_{i \geq j} (-1)^{i+j}(\sigma \circ \partial_j \circ \partial_{i-1}).
\end{align*}
Substituting $j' = j-1$ in the first sum shows the two sums cancel term by term.
\end{proof}

Thus $\{C_n(X;\mathbb{Z}), \partial_n\}$ is a \emph{chain complex}.

\begin{definition}
The \emph{$n$-th singular homology group} of $X$ is
\[
  H_n(X;\mathbb{Z})
  = \ker(\partial_n) / \operatorname{im}(\partial_{n+1})
  = Z_n / B_n,
\]
where $Z_n = \ker \partial_n$ are the \emph{cycles} and
$B_n = \operatorname{im} \partial_{n+1}$ are the \emph{boundaries}.
\end{definition}

\subsection{Basic Computations}

\begin{example}
For a point $* $: $H_0(*) = \mathbb{Z}$ and $H_n(*) = 0$ for $n > 0$.
\end{example}

\begin{example}[$n$-sphere]
\[
  H_k(S^n;\mathbb{Z}) = \begin{cases}
    \mathbb{Z} & k = 0 \text{ or } k = n, \\
    0 & \text{otherwise.}
  \end{cases}
\]
\end{example}

\begin{example}[Torus]
$H_0(T^2) = \mathbb{Z}$, $H_1(T^2) = \mathbb{Z}^2$, $H_2(T^2) = \mathbb{Z}$,
$H_k(T^2) = 0$ for $k \geq 3$.
\end{example}

\begin{example}[Real projective plane]
$H_0(\mathbb{R}P^2) = \mathbb{Z}$, $H_1(\mathbb{R}P^2) = \mathbb{Z}/2\mathbb{Z}$,
$H_k(\mathbb{R}P^2) = 0$ for $k \geq 2$.
\end{example}

\subsection{The Long Exact Sequence of a Pair}

For a pair $(X, A)$ with $A \subseteq X$, there is a long exact sequence
\[
  \cdots \to H_n(A) \xrightarrow{i_*} H_n(X)
  \xrightarrow{j_*} H_n(X, A)
  \xrightarrow{\partial} H_{n-1}(A) \to \cdots.
\]
This is a standard tool for computing homology inductively.

%% ============================================================
\section{Key Theorems: Mayer--Vietoris and Excision}
\label{sec:key}
%% ============================================================

\subsection{The Mayer--Vietoris Sequence}

\begin{theorem}[Mayer--Vietoris]
\label{thm:mayer_vietoris}
Let $X = A \cup B$ where $A$ and $B$ are open subsets (or more generally
when the interiors cover $X$). Then there is a long exact sequence
\[
  \cdots \to H_n(A \cap B)
  \xrightarrow{\Phi} H_n(A) \oplus H_n(B)
  \xrightarrow{\Psi} H_n(X)
  \xrightarrow{\Delta} H_{n-1}(A \cap B) \to \cdots,
\]
where $\Phi(\alpha) = (i_*(\alpha), -j_*(\alpha))$,
$\Psi(\alpha, \beta) = k_*(\alpha) + l_*(\beta)$,
and $\Delta$ is the connecting homomorphism.
\end{theorem}
\begin{proof}
Consider the short exact sequence of chain complexes
\[
  0 \to C_*(A \cap B) \xrightarrow{\Phi} C_*(A) \oplus C_*(B)
  \xrightarrow{\Psi} C_*(A+B) \to 0,
\]
where $C_*(A+B)$ denotes chains with support in $A$ or $B$.
The inclusion $C_*(A+B) \hookrightarrow C_*(X)$ is a quasi-isomorphism
by barycentric subdivision (one shows that every cycle in $C_*(X)$ is
homologous to one in $C_*(A+B)$ by subdividing sufficiently finely).
The snake lemma applied to the short exact sequence of chain complexes
then yields the desired long exact sequence.
\end{proof}

\begin{example}
To compute $H_n(S^n)$: Write $S^n = U \cup V$ where $U$ and $V$ are hemispheres
(both contractible), with $U \cap V \simeq S^{n-1}$.
Mayer--Vietoris gives for $n \geq 1$:
\[
  \cdots \to H_k(S^{n-1}) \to H_k(U) \oplus H_k(V) \to H_k(S^n) \to H_{k-1}(S^{n-1}) \to \cdots.
\]
Since $H_k(U) = H_k(V) = H_k(*) = 0$ for $k > 0$, this gives
$H_k(S^n) \cong H_{k-1}(S^{n-1})$ for $k > 0$ and $H_0(S^n) = \mathbb{Z}$,
from which the result follows by induction.
\end{example}

\subsection{The Excision Theorem}

\begin{theorem}[Excision]
\label{thm:excision}
Let $Z \subseteq A \subseteq X$ with $\overline{Z} \subseteq \mathrm{int}(A)$.
Then the inclusion $(X \setminus Z, A \setminus Z) \hookrightarrow (X, A)$
induces isomorphisms
\[
  H_n(X \setminus Z, A \setminus Z) \xrightarrow{\cong} H_n(X, A)
  \quad \text{for all } n.
\]
\end{theorem}
\begin{proof}
By barycentric subdivision, every singular chain in $C_n(X, A)$ can be
represented by a sum of simplices each of which maps entirely into $A$
or entirely into $X \setminus Z$. This shows that the inclusion
$C_*(X \setminus Z, A \setminus Z) \hookrightarrow C_*(X, A)$
induces a surjection on homology. Injectivity follows similarly by
subdividing null-homologous cycles.
\end{proof}

\subsection{The Euler Characteristic}

\begin{definition}
The \emph{Euler characteristic} of a finite CW-complex $X$ is
\[
  \chi(X) = \sum_{k=0}^{\dim X} (-1)^k \beta_k,
\]
where $\beta_k = \operatorname{rank}(H_k(X;\mathbb{Z}))$ are the
\emph{Betti numbers}.
\end{definition}

\begin{theorem}[Euler--Poincar\'e formula]
For a finite CW-complex with $c_k$ cells of dimension $k$:
\[
  \chi(X) = \sum_k (-1)^k c_k = \sum_k (-1)^k \beta_k.
\]
\end{theorem}

%% ============================================================
\section{Cohomology}
\label{sec:cohomology}
%% ============================================================

\subsection{Cochain Complexes}

\begin{definition}
For an abelian group $G$, the \emph{singular cochain group} is
\[
  C^n(X; G) = \mathrm{Hom}(C_n(X;\mathbb{Z}), G).
\]
The \emph{coboundary map} $\delta^n\colon C^n \to C^{n+1}$ is the
transpose of $\partial_{n+1}$: $(\delta^n f)(\sigma) = f(\partial_{n+1}\sigma)$.
\end{definition}

One has $\delta^{n+1} \circ \delta^n = 0$, giving the cochain complex
$(C^*(X;G), \delta)$.

\begin{definition}
The \emph{$n$-th cohomology group} is
\[
  H^n(X; G) = \ker \delta^n / \operatorname{im} \delta^{n-1}.
\]
\end{definition}

\subsection{The Universal Coefficient Theorem}

\begin{theorem}[Universal Coefficient Theorem for Cohomology]
\label{thm:uct}
For any topological space $X$ and abelian group $G$, there is a (non-naturally)
split short exact sequence
\[
  0 \to \mathrm{Ext}^1_\mathbb{Z}(H_{n-1}(X), G)
  \to H^n(X; G)
  \to \mathrm{Hom}(H_n(X), G)
  \to 0.
\]
\end{theorem}
\begin{proof}
The chain complex $C_*(X;\mathbb{Z})$ consists of free abelian groups. Applying
$\mathrm{Hom}(-,G)$ to the short exact sequences
$0 \to Z_n \to C_n \to B_{n-1} \to 0$ and $0 \to B_n \to Z_n \to H_n \to 0$,
and using the long exact sequences in Ext (which truncate since $Z_n, B_n$ are
free), one obtains the stated splitting.
\end{proof}

\subsection{The K\"unneth Formula}

\begin{theorem}[K\"unneth Formula]
\label{thm:kunneth}
For spaces $X$ and $Y$ with $H_*(X;\mathbb{Z})$ free abelian:
\[
  H_n(X \times Y;\mathbb{Z})
  \cong \bigoplus_{i+j=n} H_i(X;\mathbb{Z}) \otimes H_j(Y;\mathbb{Z}).
\]
For general coefficients in a field $\mathbf{k}$:
\[
  H_n(X \times Y; \mathbf{k})
  \cong \bigoplus_{i+j=n} H_i(X;\mathbf{k}) \otimes_\mathbf{k} H_j(Y;\mathbf{k}).
\]
\end{theorem}
\begin{proof}
One constructs the \emph{tensor product of chain complexes}
$(C_*(X) \otimes C_*(Y), \partial^\otimes)$ with
$\partial^\otimes(\sigma \otimes \tau) = \partial\sigma \otimes \tau
+ (-1)^{|\sigma|} \sigma \otimes \partial\tau$.
The Eilenberg--Zilber theorem provides a chain equivalence
$C_*(X \times Y) \simeq C_*(X) \otimes C_*(Y)$.
When $H_*(X)$ is free, the algebraic K\"unneth formula
$H_n(A \otimes B) \cong \bigoplus_{i+j=n} H_i(A) \otimes H_j(B)$
(for free chain complexes) completes the proof.
\end{proof}

%% ============================================================
\section{Poincar\'e Duality}
\label{sec:poincare}
%% ============================================================

\subsection{Orientation and Fundamental Class}

\begin{definition}
A connected closed orientable $n$-manifold $M$ has a unique generator
$[M] \in H_n(M;\mathbb{Z}) \cong \mathbb{Z}$, called its
\emph{fundamental class}.
\end{definition}

\subsection{Cap Product}

The \emph{cap product} $\frown\colon H_k(X;R) \otimes H^l(X;R) \to H_{k-l}(X;R)$
provides the mechanism for Poincar\'e duality.

\begin{theorem}[Poincar\'e Duality]
\label{thm:poincare_duality}
Let $M$ be a compact, connected, orientable $n$-manifold (without boundary)
with fundamental class $[M] \in H_n(M;\mathbb{Z})$. Then cap product with $[M]$
induces isomorphisms
\[
  \mathcal{D}\colon H^k(M;R) \xrightarrow{\cong} H_{n-k}(M;R)
  \quad \text{for all } k \text{ and all rings } R,
\]
given by $\mathcal{D}(\alpha) = [M] \frown \alpha$.
\end{theorem}
\begin{proof}
The proof proceeds by induction over a good cover. For an open ball $U \cong \mathbb{R}^n$,
both sides reduce to the homology of a point and the isomorphism is immediate.
For a general manifold, one covers $M$ by Euclidean charts and applies Mayer--Vietoris
inductively (the five-lemma shows the cap product map remains an isomorphism over
each step). Since $M$ is compact, the cover can be taken finite, completing
the induction.
\end{proof}

\begin{corollary}
For a compact orientable $n$-manifold $M$:
$\beta_k(M) = \beta_{n-k}(M)$ for all $k$.
In particular, $\chi(M) = 0$ if $n$ is odd.
\end{corollary}

\begin{theorem}[Poincar\'e Conjecture, dimension 3 -- Perelman 2003]
\label{thm:poincare_conj}
Every closed, simply connected 3-manifold is homeomorphic to $S^3$.
\end{theorem}
\begin{proof}[Historical note]
This was conjectured by Poincar\'e in 1904. Perelman proved it in 2002--2003 using
Hamilton's Ricci flow with surgery. The argument shows that under Ricci flow with
surgery, any simply-connected closed 3-manifold becomes extinct in finite time,
and the only closed 3-manifold that can become extinct in finite time is (homeomorphic
to) $S^3$.
\end{proof}

%% ============================================================
\section{CW-Complexes and Cellular Homology}
\label{sec:cw}
%% ============================================================

\subsection{Definition of CW-Complexes}

\begin{definition}
A \emph{CW-complex} is a topological space $X$ built inductively:
\begin{enumerate}[label=(\roman*)]
  \item $X^{-1} = \emptyset$.
  \item $X^n$ is obtained from $X^{n-1}$ by attaching $n$-cells via
    \emph{attaching maps} $\phi_\alpha\colon S^{n-1} \to X^{n-1}$:
    \[
      X^n = X^{n-1} \sqcup_{\bigsqcup_\alpha \phi_\alpha}
        \bigsqcup_\alpha D^n_\alpha.
    \]
  \item $X = \bigcup_n X^n$ with the weak topology.
\end{enumerate}
\end{definition}

\subsection{Cellular Chain Complex and Homology}

\begin{definition}
The \emph{cellular chain group} $C_n^{\mathrm{CW}}(X)$ is the free abelian group
generated by the $n$-cells of $X$.
The cellular boundary map $d_n\colon C_n^{\mathrm{CW}} \to C_{n-1}^{\mathrm{CW}}$
is given by the \emph{degree} of attaching maps:
\[
  d_n(e^n_\alpha) = \sum_\beta d_{\alpha\beta} e^{n-1}_\beta,
\]
where $d_{\alpha\beta} \in \mathbb{Z}$ is the degree of the composite
$S^{n-1} \xrightarrow{\phi_\alpha} X^{n-1} \twoheadrightarrow S^{n-1}_\beta$.
\end{definition}

\begin{theorem}[Cellular homology equals singular homology]
$H_n(C_*^{\mathrm{CW}}(X), d) \cong H_n(X;\mathbb{Z})$ for all $n$.
\end{theorem}

\begin{example}
CW-structures of standard spaces:
\begin{center}
\begin{tabular}{lll}
\toprule
Space & Cell structure & Homology \\
\midrule
$S^n$ & one $0$-cell, one $n$-cell & $H_0 = H_n = \mathbb{Z}$, rest $0$ \\
$\mathbb{R}P^n$ & one cell per dimension $0,\ldots,n$ & alternating $\mathbb{Z}$ and $\mathbb{Z}/2$ \\
$T^2$ & one $0$-, two $1$-, one $2$-cell & $H_0=\mathbb{Z}$, $H_1=\mathbb{Z}^2$, $H_2=\mathbb{Z}$ \\
$\mathbb{R}P^2$ & one $0$-, one $1$-, one $2$-cell & $H_0=\mathbb{Z}$, $H_1=\mathbb{Z}/2$, $H_2=0$ \\
\bottomrule
\end{tabular}
\end{center}
\end{example}

%% ============================================================
\section{Higher Homotopy Groups}
\label{sec:higher}
%% ============================================================

\subsection{Definition}

\begin{definition}
The \emph{$n$-th homotopy group} of $(X, x_0)$ is
\[
  \pi_n(X, x_0) = [(I^n, \partial I^n), (X, x_0)],
\]
the set of homotopy classes of maps $(I^n, \partial I^n) \to (X, x_0)$
(relative to the boundary).
\end{definition}

For $n \geq 2$, $\pi_n(X, x_0)$ is abelian. For $n = 1$, this recovers $\pi_1$.

\begin{example}
$\pi_n(S^n) \cong \mathbb{Z}$ for all $n \geq 1$ (generator: identity map).

$\pi_3(S^2) \cong \mathbb{Z}$ (generated by the Hopf fibration $S^3 \to S^2$).
\end{example}

\subsection{The Hurewicz Theorem}

\begin{theorem}[Hurewicz Isomorphism]
\label{thm:hurewicz}
Let $X$ be an $(n-1)$-connected space (i.e., $\pi_k(X) = 0$ for $0 \leq k < n$)
with $n \geq 2$. Then $H_k(X) = 0$ for $0 < k < n$ and the
\emph{Hurewicz homomorphism}
\[
  h\colon \pi_n(X, x_0) \to H_n(X;\mathbb{Z}),
  \quad [f] \mapsto f_*(\iota_n),
\]
(where $\iota_n \in H_n(S^n) \cong \mathbb{Z}$ is the fundamental class)
is an isomorphism.
\end{theorem}
\begin{proof}
The proof uses the Serre spectral sequence applied to the path-loop fibration
$\Omega X \to PX \to X$ (where $PX$ is contractible). Since $X$ is $(n-1)$-connected,
$H_k(\Omega X) = 0$ for $0 < k < n-1$, and the spectral sequence collapses in
the relevant range to give the isomorphism. For $n=1$ (simply connected case),
a direct simplicial argument suffices: every 1-cycle defines a loop, and this
defines the Hurewicz map as an isomorphism.
\end{proof}

\begin{theorem}[Whitehead's Theorem]
Let $f\colon X \to Y$ be a map between simply connected CW-complexes.
If $f_*\colon H_n(X;\mathbb{Z}) \xrightarrow{\cong} H_n(Y;\mathbb{Z})$
for all $n$, then $f$ is a homotopy equivalence.
\end{theorem}

\begin{theorem}[Adams -- Hopf Invariant 1]
\label{thm:adams}
There exist maps $f\colon S^{2n-1} \to S^n$ of Hopf invariant $1$ only for
$n = 1, 2, 4, 8$.
\end{theorem}
\begin{proof}[Historical note]
Proved by J. F. Adams in 1960 using secondary cohomology operations (the
Adams operations in K-theory also give an elegant proof). The result implies
that $\mathbb{R}^n$ admits a normed division algebra structure only for
$n = 1, 2, 4, 8$ (i.e., $\mathbb{R}, \mathbb{C}, \mathbb{H}, \mathbb{O}$).
\end{proof}

%% ============================================================
\section{Applications}
\label{sec:applications}
%% ============================================================

\subsection{The Brouwer Fixed-Point Theorem}

\begin{theorem}[Brouwer Fixed-Point Theorem]
\label{thm:brouwer}
Every continuous map $f\colon D^n \to D^n$ has a fixed point.
\end{theorem}
\begin{proof}
Suppose $f$ has no fixed point. Define a retraction
$r\colon D^n \to S^{n-1}$ by: for each $x \in D^n$, let $r(x)$ be the point
where the ray from $f(x)$ through $x$ meets $S^{n-1}$. This map is continuous.
But $r \circ \iota = \mathrm{id}_{S^{n-1}}$ where $\iota\colon S^{n-1} \hookrightarrow D^n$
is the inclusion. Applying $H_{n-1}$:
\[
  \mathbb{Z} \cong H_{n-1}(S^{n-1})
  \xrightarrow{r_*} H_{n-1}(D^n) = 0
  \xrightarrow{\iota_*} H_{n-1}(S^{n-1}) \cong \mathbb{Z},
\]
with $r_* \circ \iota_* = \mathrm{id}$. This requires $\mathbb{Z}$ to
inject into $0$, a contradiction.
\end{proof}

\subsection{The Borsuk--Ulam Theorem}

\begin{theorem}[Borsuk--Ulam]
\label{thm:borsuk_ulam}
For every continuous map $f\colon S^n \to \mathbb{R}^n$, there exists a point
$x \in S^n$ with $f(x) = f(-x)$.
\end{theorem}
\begin{proof}
Assume $f(x) \neq f(-x)$ for all $x$. Define
$g\colon S^n \to S^{n-1}$ by $g(x) = (f(x) - f(-x)) / \|f(x) - f(-x)\|$.
Then $g(-x) = -g(x)$, so $g$ is an \emph{odd} map.
But an odd map $S^n \to S^{n-1}$ has odd degree as a map $S^{n-1} \to S^{n-1}$
when restricted to the equator -- in particular, its restriction induces a
non-trivial map on $H_{n-1}$. Since any odd map $S^n \to S^{n-1}$ must factor
through an $(n-1)$-dimensional space but $H_n(S^n) \neq 0$, one obtains a
contradiction from the Lefschetz--Hopf theory applied to the degree.
\end{proof}

\begin{corollary}
At any moment, there exist two antipodal points on Earth's surface with identical
temperature and pressure.
\end{corollary}

\subsection{The Hairy Ball Theorem}

\begin{theorem}[Hairy Ball Theorem]
\label{thm:hairy_ball}
There is no non-vanishing continuous tangent vector field on $S^{2n}$
(the even-dimensional sphere).
\end{theorem}
\begin{proof}
Suppose $v\colon S^{2n} \to \mathbb{R}^{2n+1}$ is a non-vanishing tangent
vector field (so $v(x) \perp x$ and $v(x) \neq 0$ for all $x \in S^{2n}$).
Normalise to $\|v(x)\| = 1$. Define a homotopy
$H_t(x) = \cos(\pi t) x + \sin(\pi t) v(x)$, $t \in [0,1]$.
Then $H_0 = \mathrm{id}_{S^{2n}}$ and $H_1 = -\mathrm{id}_{S^{2n}}$ (the antipodal map).
The antipodal map on $S^{2n}$ has degree $(-1)^{2n+1} = -1$.
But since $H_0$ and $H_1$ are homotopic, they have the same degree,
so $\deg(\mathrm{id}) = 1 = -1 = \deg(-\mathrm{id})$, a contradiction.
\end{proof}

\begin{remark}
On $S^{2n-1}$ (odd-dimensional spheres), non-vanishing vector fields do exist:
the unit tangent vector field $v(x_1,\ldots,x_{2n}) = (-x_2,x_1,\ldots,-x_{2n},x_{2n-1})$
provides one.
\end{remark}

\subsection{Surface Classification}

\begin{theorem}[Classification of compact surfaces]
Every compact connected surface is homeomorphic to exactly one of:
\begin{enumerate}[label=(\roman*)]
  \item $\Sigma_g$ (orientable, genus $g \geq 0$): $\chi = 2-2g$.
  \item $N_k$ (non-orientable, $k \geq 1$ crosscaps): $\chi = 2-k$.
\end{enumerate}
The Euler characteristic and orientability are complete invariants.
\end{theorem}

%% ============================================================
\section*{References}
%% ============================================================

\begin{thebibliography}{99}

\bibitem{Hatcher}
A.~Hatcher.
\newblock \textit{Algebraic Topology}.
\newblock Cambridge University Press, Cambridge, 2002.
\newblock Free online: \url{https://pi.math.cornell.edu/~hatcher/AT/ATpage.html}.

\bibitem{Bredon}
G.~E.~Bredon.
\newblock \textit{Topology and Geometry}.
\newblock Graduate Texts in Mathematics 139, Springer, New York, 1993.

\bibitem{May}
J.~P.~May.
\newblock \textit{A Concise Course in Algebraic Topology}.
\newblock University of Chicago Press, Chicago, 1999.

\bibitem{Munkres}
J.~R.~Munkres.
\newblock \textit{Elements of Algebraic Topology}.
\newblock Addison-Wesley, Menlo Park, CA, 1984.

\bibitem{Spanier}
E.~H.~Spanier.
\newblock \textit{Algebraic Topology}.
\newblock Springer, New York, 1981.

\bibitem{Adams}
J.~F.~Adams.
\newblock On the non-existence of elements of Hopf invariant one.
\newblock \textit{Annals of Mathematics}, 72:20--104, 1960.

\bibitem{Perelman1}
G.~Perelman.
\newblock The entropy formula for the Ricci flow and its geometric applications.
\newblock \textit{arXiv:math/0211159}, 2002.

\bibitem{Perelman2}
G.~Perelman.
\newblock Ricci flow with surgery on three-manifolds.
\newblock \textit{arXiv:math/0303109}, 2003.

\bibitem{Milnor}
J.~Milnor.
\newblock \textit{Morse Theory}.
\newblock Princeton University Press, Princeton, 1963.

\end{thebibliography}

\end{document}
